\documentclass[letterpaper]{article}

\usepackage[letterpaper,margin=1in]{geometry}

% Compile from the repository root so that this path resolves to
% the pinned WSP submodule.
\input{wsp/testing/test-case-library.tex}

\begin{document}

\def\TCID{%
  TC-0016
}

\def\TCTitle{%
  Persisted state in a form file
}

\def\TCPurpose{%
  Verify that the property groups VB6 writes for the common controls are read
  from a form file, and that the binary stream of the TreeView is reported and
  not read.
}

\def\TCPriority{%
  \TCPriorityHigh
}

\def\TCRequirementRef{%
  REQ-0016, REQ-0023
}

\def\TCDesignRef{%
  \texttt{ADR-0006} and \texttt{ADR-0010} in WFC
}

\def\TCFeatureRef{%
  \texttt{tests/corpus/194-comctl-persisted.frm} and
  \texttt{tests/corpus/194-comctl-persisted.vbp}; the form code and the
  controls of \texttt{include/wfc/mscomctl/mscomctl.inc} and the runner
  \texttt{tests/tc-0016-comctl-persisted.ps1}
}

\def\TCTestTechnique{%
  Scenario-based testing of one form: each Immediate-window line is compared
  with the recorded outcome (ISO/IEC/IEEE 29119-4 experience-based and
  specification-based techniques, with expected output comparison).
}

\def\TCPreconditions{%
  A \texttt{wfc} of the same version as the library (WFC, the interpreter and
  compiler these controls belong to, runs the case: the library holds the
  program, WFC the runner). The files of the case are in
  \texttt{tests/corpus}: \texttt{194-comctl-persisted.vbp},
  \texttt{194-comctl-persisted.frm}, \texttt{194-comctl-persisted.expected},
  \texttt{194-comctl-persisted.debug}, \texttt{194-comctl-persisted.ui},
  \texttt{194-comctl-persisted.frx}, \texttt{194-comctl-persisted.red.bmp}.
}

\def\TCEnvironment{%
  Windows with \texttt{wfc.exe} of the x64 Debug build; the form runs in WFC's
  headless host (no desktop is needed). WFC registers the case with CTest and
  runs \texttt{wfc --interpret 194-comctl-persisted.vbp}.
}

\def\TCAssumptions{%
  The scripted input \texttt{194-comctl-persisted.ui} is the only input; the
  case does not read the registry or other files except the picture named in
  the form.
}

\def\TCInputData{%
  A form file with \texttt{Tabs}, \texttt{Panels}, \texttt{Buttons} with
  \texttt{ButtonMenus}, \texttt{ColumnHeader} and \texttt{Images} groups, an
  \texttt{Object =} line and an \texttt{OleObjectBlob} for the TreeView; the
  picture in \texttt{194-comctl-persisted.frx}.
}

\def\TCInitialState{%
  A fresh \texttt{wfc} process; the program creates its own controls from the
  form file and starts with no state from a previous run.
}

\def\TCProcedure{%
  \begin{enumerate}
    \item Load the form and print the tabs (count, index of \texttt{adv}, tool
      tip, selected key).
    \item Print the panels (count, text, AutoSize, Width, style).
    \item Print the buttons, the style of \texttt{open} and the text of its
      menu.
    \item Print the column headers, the image list and the TreeView node
      count.
  \end{enumerate}
}

\def\TCExpectedResult{%
  \begin{enumerate}
    \item Two tabs, \texttt{adv} at index 2, the tool tip \texttt{tip} and the
      first tab selected (\texttt{gen}).
    \item Two panels, \texttt{Ready}, AutoSize 1, width 4000 and style 6.
    \item Two buttons, style 5 and the menu text \texttt{Recent}.
    \item Two columns (\texttt{Size} at 2000), one image with key \texttt{red}
      at index 1 and width 32.
    \item The TreeView has no nodes (0) and \texttt{WFC0293} is written to the
      Immediate output.
    \item \texttt{wfc} exits with code 0, writes the exact text of
      \texttt{194-comctl-persisted.expected} to standard output, and writes
      the exact Immediate-window text of \texttt{194-comctl-persisted.debug}
      to standard error. Any difference fails the case.
  \end{enumerate}
}

\def\TCPostConditions{%
  The form is closed by the last \texttt{close Form1} line of the script;
  nothing is left running and no file is written.
}

\TCTemplate

\end{document}
